% Fragment partage : inclus par trajets-confiance.tex (dossier autonome)
% et par dossier-conception-alanya.tex (dossier client assemble).
% Ne pas y mettre de preambule ni de \part : c'est l'appelant qui les pose.

\chapter{La propriété de sûreté}

\section{Le désaccord à trancher en premier}

Une fonctionnalité de suivi de trajet peut se construire autour de deux promesses
différentes, et elles n'ont pas le même coût ni la même solidité~:

\begin{enumerate}[leftmargin=1.6em]
  \item \textbf{« vos proches voient où vous êtes »} --- la promesse est la \emph{trace}~;
  \item \textbf{« si vous ne confirmez pas, vos proches sont prévenus »} --- la promesse est
        l'\emph{échéance}.
\end{enumerate}

La première dépend du GPS, de la batterie, du réseau, du système d'exploitation et du bon
vouloir d'iOS en arrière-plan. Elle tombe en panne régulièrement, et une fonctionnalité de
sûreté qui tombe en panne régulièrement est pire qu'absente~: on lui fait confiance.

La seconde ne dépend que d'une ligne dans \code{job\_queue} et d'une horloge serveur.

\begin{note}[Le postulat du volet]
\textbf{La garantie de sûreté, c'est l'échéance --- pas la trace.} La trace est un confort~:
elle rassure pendant, et elle aide à retrouver quelqu'un après. L'échéance est le contrat.
Tout le reste du dossier découle de là.
\end{note}

\section{Trois conséquences, à ne jamais perdre de vue}

\textbf{Conséquence 1 --- perdre la position n'est pas un incident.} Un téléphone dans un
tunnel, un GPS coupé, une application tuée par le système~: ce sont des \emph{informations},
affichées honnêtement (« dernier point il y a 4 minutes »), et elles ne déclenchent jamais
d'alerte à elles seules. Confondre « je ne vous vois plus » et « il vous est arrivé quelque
chose » est le plus court chemin vers un cercle qui n'ouvre plus les notifications.

\textbf{Conséquence 2 --- iOS n'a pas à égaler Android.} On ne peut pas garantir un flux de
positions continu sur iOS en arrière-plan, et il est inutile de le promettre. Le filet tient
quand même, parce qu'il est côté serveur. Le dossier dit ce qui est garanti et ce qui ne
l'est pas, plutôt que d'annoncer une parité qui ne sera pas tenue.

\textbf{Conséquence 3 --- on ne clôt jamais automatiquement sur une arrivée détectée.} Le
GPS \emph{propose}, l'humain \emph{dispose}. Une détection d'arrivée pose la question~; seule
la personne y répond. Le coût d'un faux positif devient une question à balayer, jamais un
filet rompu.

\begin{alerte}[Le filet dépend d'un drapeau d'environnement]
\code{startJobWorker()} est conditionné par \code{BROADCAST\_\allowbreak WORKER\_\allowbreak ENABLED !== 'true'}
(\code{src/services/jobQueue.js}). Aujourd'hui, si le drapeau est absent, la file ne tourne
pas et seuls les diffusions et le rattrapage de bienvenue en pâtissent --- fâcheux, pas
grave. Avec les trajets, \textbf{aucune échéance ne se déclenche plus~: on promet un filet
qui n'existe pas}. Deux mesures dans le même lot~: renommer le drapeau en
\code{JOB\_WORKER\_ENABLED}, et refuser de servir \mPOST{}~\code{/api/trips} si le worker
n'a pas démarré. Un refus explicite vaut mieux qu'une promesse silencieuse.
\end{alerte}

\chapter{La machine à états}

\section{Les huit états}

\begin{center}
\begin{tabularx}{\linewidth}{@{}p{3.9cm}p{1.7cm}X@{}}
\toprule
\thd{État} & \thd{Nature} & \thd{Sens} \\
\midrule
\code{active}            & ouvert   & Partage en cours, les positions circulent. \\
\code{awaiting\_confirm} & ouvert   & La question est posée. Le partage continue. \\
\code{alert}             & ouvert   & Grâce écoulée sans réponse. Le cercle est prévenu. \\
\code{sos}               & ouvert   & Alerte explicite. Strictement plus fort que
                                      \code{alert}. \\
\code{closed\_confirmed} & terminal & La personne a confirmé son arrivée. \\
\code{closed\_cancelled} & terminal & Arrêté par le propriétaire avant toute échéance. \\
\code{closed\_expired}   & terminal & Plafond dur atteint après une alerte. \\
\code{closed\_unwatched} & terminal & Plus aucun destinataire. \\
\bottomrule
\end{tabularx}
\end{center}

\begin{note}[Pourquoi des libellés anglais dans le schéma]
Le reste de la base fait de même --- \code{platform ENUM('android','ios','web','unknown')}
(\code{migrations/\allowbreak 020\_\allowbreak user\_\allowbreak push\_\allowbreak devices.sql}), \code{kind = 'trust'}
(\code{migrations/\allowbreak 050\_\allowbreak contact\_\allowbreak list\_\allowbreak kind.sql}). Les libellés français sont l'affaire de la
couche de présentation et de \code{lib/l10n/}, jamais celle du schéma.
\end{note}

\begin{alerte}[\code{ENUM} plutôt qu'un entier --- et le contre-exemple est dans la base]
À taille de stockage \textbf{identique} --- un octet dans les deux cas --- l'entier ne
survit pas à la relecture. \code{message.type} est un \code{SMALLINT} documenté par un
commentaire SQL~: \code{COMMENT '0=text 1=image 2=video 3=audio 4=file 5=location'}
(\code{migrations/\allowbreak 001\_\allowbreak initial\_\allowbreak schema.sql}). Ce
commentaire s'arrête à~5, alors que les types \textbf{6, 7 et 8 existent} --- message
système, carte de contact, appel à l'action de bienvenue. Le commentaire a pourri~; la base
ne sait plus ce qu'elle contient, et un dump de production ne se lit plus sans le code sous
les yeux.

Un \code{ENUM} ne peut pas pourrir~: la liste des valeurs \emph{est} le schéma, elle sort
avec \code{SHOW CREATE TABLE}, et le serveur refuse ce qui n'y figure pas. C'est déjà le
choix retenu partout où la base nomme un état~: \code{appState}, \code{block\_\allowbreak
type}, \code{login\_\allowbreak method}, \code{previewMode}, et le \code{status} des
diffusions.
\end{alerte}

\begin{note}[La règle appliquée ici, et ses trois pièges]
\textbf{Ensemble fermé, qui ne change que par décision de conception} $\rightarrow$
\code{ENUM}~: \code{trip.state} --- c'est une machine à états, y ajouter un état \emph{doit}
coûter une migration, c'est une garantie et non une gêne --- ainsi que \code{trip.kind} et
\code{trip\_\allowbreak watcher.state}.

\textbf{Ensemble ouvert, écrit depuis de nombreux endroits} $\rightarrow$ \code{VARCHAR}~:
\code{trip\_\allowbreak event.kind}, qui gagnera des valeurs au fil des versions, comme
\code{job\_\allowbreak queue.kind} avant lui.

\textbf{Entier} $\rightarrow$ seulement si la valeur est \emph{ordinale} et qu'on la compare
avec \code{<} ou \code{>}. Aucune colonne de ce volet n'est dans ce cas.

Les trois pièges de l'\code{ENUM}, à connaître avant d'écrire la migration~:
\begin{enumerate}[leftmargin=1.6em]
  \item \code{WHERE state = 1} ne compare pas au texte «~1~»~: il désigne la
        \textbf{première valeur déclarée}. On compare toujours à une chaîne.
  \item \code{ORDER BY state} trie par ordre de \emph{déclaration}, pas alphabétiquement.
        On déclare donc les valeurs dans l'ordre du cycle de vie --- ici, c'est un avantage.
  \item Ajouter une valeur \textbf{en fin de liste} se fait sans reconstruire la table, tant
        qu'on reste sous 255 valeurs et donc sur un octet. L'insérer au milieu, la renommer
        ou réordonner force au contraire une reconstruction complète. On ajoute toujours à
        la fin.
\end{enumerate}
\end{note}

\section{L'automate}

\begin{center}
\begin{tikzpicture}[
  x=1cm, y=1cm,
  e/.style={draw=AlanyaLine, fill=AlanyaSurface, rounded corners=3pt,
            font=\sffamily\small, inner sep=6pt, minimum height=0.85cm, text=AlanyaInk},
  ok/.style={e, draw=AlanyaSuccess, line width=0.9pt, fill=AlanyaSuccess!8},
  ko/.style={e, draw=danger, line width=0.9pt, fill=danger!6},
  fin/.style={e, draw=AlanyaIndigo, fill=AlanyaContainer, text=AlanyaIndigoDark},
  a/.style={-{Stealth[length=2.2mm]}, draw=AlanyaGray, semithick},
  l/.style={font=\sffamily\scriptsize, text=AlanyaGray, fill=white, inner sep=2pt}
]
  \node[e]   (act) at (0,0)       {\code{active}};
  \node[e]   (awa) at (4.9,0)     {\code{awaiting\_confirm}};
  \node[ok]  (con) at (10.4,0)    {\code{closed\_confirmed}};
  \node[ko]  (ale) at (4.9,-2.6)  {\code{alert}};
  \node[ko]  (sos) at (0,-2.6)    {\code{sos}};
  \node[fin] (exp) at (10.4,-2.6) {\code{closed\_expired}};

  \draw[a] (act) -- node[l, above] {ETA, ou arrivée tenue 60~s} (awa);
  \draw[a] (awa) -- node[l, above] {confirmation}               (con);
  \draw[a] (ale) -- node[l, above right, pos=0.15, inner sep=1pt]
           {confirmation tardive} (con.south west);
  \draw[a] (awa) -- node[l, right] {grâce écoulée}              (ale);
  \draw[a] (act) -- node[l, left]  {bouton SOS}                 (sos);
  \draw[a] (awa.south west) -- node[l, pos=0.5, below left, inner sep=1pt]
           {bouton SOS} (sos.north east);
  \draw[a] (ale) -- node[l, above] {plafond dur}                (exp);

  % Prolongation : boucle par le haut, elle ne doit croiser aucune transition.
  \draw[a] (awa.north) -- ++(0,0.85) -- node[l, above] {prolongation}
           ++(-4.9,0) -- (act.north);

  % Arret volontaire : route par le bas, sous tous les etats.
  \draw[a] (act.west) -- (-1.95,0) -- (-1.95,-4.3)
           -- node[l, below, pos=0.60] {arrêt, ou plus aucun destinataire}
           (10.4,-4.3) -- (exp.south);
\end{tikzpicture}
\end{center}

Les deux états terminaux \code{closed\_cancelled} et \code{closed\_unwatched} partagent la
même arête de sortie sur le schéma~; ils diffèrent par leur \code{close\_reason} et par ce
qui est notifié au cercle.

\section{Les transitions}

\begin{center}
\begin{tabularx}{\linewidth}{@{}p{3.0cm}p{3.0cm}p{2.2cm}X@{}}
\toprule
\thd{De} & \thd{Vers} & \thd{Qui} & \thd{Déclencheur} \\
\midrule
---                 & \code{active}   & client & \mPOST{}~\code{/api/trips} \\
---                 & \code{sos}      & client & \mPOST{}~\code{/api/trips/sos},
                                                 aucun trajet ouvert \\
\code{active}       & \code{awaiting\_confirm} & job / serveur &
  \code{trip\_eta\_due}, ou rayon tenu 60~s à l'ingestion d'un point \\
\code{awaiting\_confirm} & \code{active} & client & prolongation \\
\code{awaiting\_confirm} & \code{closed\_confirmed} & client & confirmation \\
\code{awaiting\_confirm} & \code{alert} & job & \code{trip\_grace\_expired} \\
\code{alert}        & \code{closed\_confirmed} & client & confirmation tardive \\
ouvert              & \code{sos}      & client & bouton SOS après décompte \\
\code{alert} / \code{sos} & \code{closed\_expired} & job &
  \code{trip\_max\_duration} \\
\code{active} / \code{awaiting\_confirm} & \code{closed\_cancelled} & client & arrêt \\
tout ouvert         & \code{closed\_unwatched} & job & \code{trip\_no\_watcher} \\
\bottomrule
\end{tabularx}
\end{center}

\section{Quatre invariants}

\begin{enumerate}[leftmargin=1.6em]
  \item \textbf{Seul le propriétaire clôt.} Un destinataire ne peut qu'acquitter
        (\mPOST{}~\code{/api/trips/:id/seen}), ce qui s'inscrit dans \code{trip\_event} et
        remonte au propriétaire (« Maman a vu »). Personne d'autre ne décide qu'une personne
        est arrivée.
  \item \textbf{Une alerte émise ne se rétracte jamais.} Confirmer après une alerte émet un
        événement de résolution~; cela n'efface pas l'alerte, ni chez les destinataires, ni
        dans le journal.
  \item \textbf{Terminal est terminal.} Repartir, c'est un nouveau trajet.
  \item \textbf{Un seul trajet ouvert par utilisateur} --- 409 \code{TRIP\_ALREADY\_ACTIVE}.
\end{enumerate}

\section{Les cas limites}

\begin{center}
\begin{tabularx}{\linewidth}{@{}p{4.2cm}X@{}}
\toprule
\thd{Situation} & \thd{Comportement} \\
\midrule
\thd{Application tuée} & Le service en avant-plan tient le flux. Si le système tue quand
même, le serveur constate la péremption ($3\times$ battement $+$ 30~s), pose
\code{stale = 1} et émet \code{trip:stale}. \textbf{Il n'alerte pas.} \\
\thd{GPS coupé, permission retirée} & Le client émet \code{trip:signal} avec son motif~;
le cercle lit « position indisponible » et le dernier point daté. \textbf{La chaîne
d'échéance continue}~: perdre le signal n'exempte pas de confirmer. \\
\thd{Réseau coupé} & Tampon local borné en Drift, vidé à la reconnexion avec les
horodatages réels. Même chorégraphie de reprise que \code{RESUME\_ACK\_MS}
(\code{src/\allowbreak socket/\allowbreak state/\allowbreak callState.js}). \\
\thd{Deux appareils du propriétaire} & Exactement le problème des appels. Un lien
\code{(tripId, ownerId) $\rightarrow$ deviceId} sur le modèle de
\code{src/\allowbreak socket/\allowbreak state/\allowbreak callDeviceOwnership.js}. Sans lui, deux appareils entrelacent leurs
positions et la trace zigzague. \\
\thd{Trajet plus long que prévu} & Prolongations illimitées, plus un plafond dur de 12~h qui
force \code{awaiting\_confirm}. Un trajet n'est pas un traceur permanent. \\
\thd{Compte supprimé ou exclu} & Clôture en \code{closed\_cancelled}, destinataires
notifiés. \\
\thd{Horloge du client fausse} & Toute l'arithmétique d'échéance est serveur, en UTC, via
\code{job\_queue.run\_after}. Le client n'envoie qu'une durée ou une heure~; c'est
\code{eta\_at} calculé par le serveur qui fait foi. \\
\bottomrule
\end{tabularx}
\end{center}

\chapter{Modèle de données}

\section{Quatre tables}

\begin{lstlisting}[language=SQL, caption={\code{migrations/051\_trips.sql} --- extrait}]
-- 051_trips.sql -- trajets de confiance : suivi temps reel, echeance et SOS.
-- L'audience est figee au demarrage (instantane de la liste Confiance),
-- sur le modele de broadcast_delivery (040_broadcast.sql).

CREATE TABLE IF NOT EXISTS trip (
  id            BIGINT       NOT NULL AUTO_INCREMENT,
  owner_id      INT          NOT NULL,
  client_id     VARCHAR(64)  NOT NULL,          -- creation idempotente
  kind          ENUM('taxi','meeting','sos')      NOT NULL DEFAULT 'meeting',
  state         ENUM('active','awaiting_confirm','alert','sos',
                     'closed_confirmed','closed_cancelled',
                     'closed_expired','closed_unwatched')
                                                 NOT NULL DEFAULT 'active',
  dest_lat      DECIMAL(9,6) NULL,
  dest_lng      DECIMAL(9,6) NULL,
  dest_label    VARCHAR(160) NULL,               -- geocode UNE fois a la creation
  dest_radius_m SMALLINT     NOT NULL DEFAULT 100,
  eta_at        DATETIME     NULL,
  grace_minutes SMALLINT     NOT NULL DEFAULT 10,
  extensions    SMALLINT     NOT NULL DEFAULT 0,
  note          VARCHAR(200) NULL,
  owner_device  VARCHAR(128) NULL,               -- seul appareil emetteur
  last_lat      DECIMAL(9,6) NULL,
  last_lng      DECIMAL(9,6) NULL,
  last_acc_m    SMALLINT     NULL,
  last_battery  TINYINT      NULL,
  last_at       DATETIME     NULL,
  stale         TINYINT(1)   NOT NULL DEFAULT 0,
  started_at    DATETIME     NOT NULL DEFAULT CURRENT_TIMESTAMP,
  alerted_at    DATETIME     NULL,
  closed_at     DATETIME     NULL,
  close_reason  VARCHAR(30)  NULL,
  points_purged_at DATETIME  NULL,
  PRIMARY KEY (id),
  UNIQUE KEY uq_trip_client (owner_id, client_id),
  KEY idx_trip_owner  (owner_id, started_at),
  KEY idx_trip_open   (state, eta_at),
  KEY idx_trip_purge  (closed_at, points_purged_at),
  CONSTRAINT fk_trip_owner FOREIGN KEY (owner_id)
    REFERENCES users(alanyaID) ON DELETE CASCADE
) ENGINE=InnoDB DEFAULT CHARSET=utf8mb4 COLLATE=utf8mb4_unicode_ci;

CREATE TABLE IF NOT EXISTS trip_watcher (
  trip_id       BIGINT   NOT NULL,
  alanyaID      INT      NOT NULL,
  state         ENUM('active','revoked','left') NOT NULL DEFAULT 'active',
  revoke_reason VARCHAR(20) NULL,                -- blocked | removed | left
  msgID         INT      NULL,                   -- la carte type 9 chez ce membre
  conversID     INT      NULL,
  added_at      DATETIME NOT NULL DEFAULT CURRENT_TIMESTAMP,
  seen_at       DATETIME NULL,
  revoked_at    DATETIME NULL,
  PRIMARY KEY (trip_id, alanyaID),
  KEY idx_tw_user (alanyaID, trip_id),
  CONSTRAINT fk_tw_trip FOREIGN KEY (trip_id)
    REFERENCES trip(id) ON DELETE CASCADE,
  CONSTRAINT fk_tw_user FOREIGN KEY (alanyaID)
    REFERENCES users(alanyaID) ON DELETE CASCADE
) ENGINE=InnoDB DEFAULT CHARSET=utf8mb4 COLLATE=utf8mb4_unicode_ci;

CREATE TABLE IF NOT EXISTS trip_point (
  id          BIGINT       NOT NULL AUTO_INCREMENT,
  trip_id     BIGINT       NOT NULL,
  lat         DECIMAL(9,6) NOT NULL,
  lng         DECIMAL(9,6) NOT NULL,
  acc_m       SMALLINT     NULL,
  speed_kmh   SMALLINT     NULL,
  battery     TINYINT      NULL,
  recorded_at DATETIME(3)  NOT NULL,             -- horloge client
  received_at DATETIME     NOT NULL DEFAULT CURRENT_TIMESTAMP,
  PRIMARY KEY (id),
  KEY idx_point_trip (trip_id, recorded_at),
  CONSTRAINT fk_point_trip FOREIGN KEY (trip_id)
    REFERENCES trip(id) ON DELETE CASCADE
) ENGINE=InnoDB DEFAULT CHARSET=utf8mb4 COLLATE=utf8mb4_unicode_ci;

CREATE TABLE IF NOT EXISTS trip_event (
  id         BIGINT      NOT NULL AUTO_INCREMENT,
  trip_id    BIGINT      NOT NULL,
  kind       VARCHAR(24) NOT NULL,
  actor_id   INT         NULL,                   -- NULL = systeme
  meta       JSON        NULL,
  created_at DATETIME    NOT NULL DEFAULT CURRENT_TIMESTAMP,
  PRIMARY KEY (id),
  KEY idx_event_trip (trip_id, id),
  CONSTRAINT fk_event_trip FOREIGN KEY (trip_id)
    REFERENCES trip(id) ON DELETE CASCADE
) ENGINE=InnoDB DEFAULT CHARSET=utf8mb4 COLLATE=utf8mb4_unicode_ci;

INSERT IGNORE INTO scheduler_leases (name) VALUES ('trip_nightly_purge');
\end{lstlisting}

Les valeurs de \code{trip\_event.kind}~: \code{started}, \code{extended},
\code{arrival\_detected}, \code{eta\_due}, \code{confirmed}, \code{alerted}, \code{sos},
\code{resolved}, \code{watcher\_seen}, \code{watcher\_revoked}, \code{signal\_lost},
\code{signal\_back}, \code{low\_battery}, \code{device\_takeover}, \code{closed}.

\begin{alerte}[Les migrations sont appliquées à la main]
Comme le rappelle déjà l'en-tête de \code{migrations/\allowbreak 038\_contact\_lists.sql}, il n'existe ni
exécuteur ni table de suivi. \code{051\_trips.sql} doit être appliqué \textbf{avant} de
déployer le code qui nomme ces tables. La base est distante~: prévoir la fenêtre.
\end{alerte}

\section{Pourquoi \code{DECIMAL(9,6)} et aucun index spatial}

Six décimales valent environ 11~cm --- très au-delà de ce qu'un GPS de téléphone sait
produire. Le type est exact, comparable, se sauvegarde en dump texte sans surprise, et il
rend facile de dire honnêtement quelle précision on stocke.

Un type \code{POINT} et un index \code{SPATIAL} n'apporteraient rien~: on ne cherche
\textbf{jamais} de proximité entre trajets. L'arrivée s'évalue par trajet, contre une
destination unique, en mémoire, à l'ingestion du point. Les cinquante migrations existantes
n'ont aucun index spatial~; en introduire un pour un calcul qu'on ne fait pas serait un coût
d'exploitation gratuit.

\section{L'instantané de l'audience}

\code{trip\_watcher} est peuplé \textbf{une fois}, au démarrage, à partir de \emph{tous} les
membres de la liste \code{kind = 'trust'} du propriétaire. C'est le même geste que
\code{broadcast\_delivery} (\code{migrations/\allowbreak 040\_broadcast.sql})~: matérialiser l'audience
pour pouvoir dire, après coup, qui a été destinataire.

Trois raisons~:

\begin{enumerate}[leftmargin=1.6em]
  \item un trajet est une pièce d'archive --- si l'alerte part, il faut pouvoir répondre à
        « qui a été prévenu~?~»~;
  \item les modifications de \code{contact\_list\_member} ne doivent pas se propager à un
        trajet en cours, sinon l'audience devient indécidable~;
  \item cela neutralise la faille décrite plus bas~: supprimer sa liste Confiance pendant un
        trajet n'ampute pas le trajet en cours.
\end{enumerate}

\begin{note}[Pas de sous-sélection --- le cercle \emph{est} l'audience]
Le composeur ne propose aucune case à cocher~: démarrer un trajet notifie les cinq membres
du cercle, en bloc. C'est ce qui permet de partir en un appui, et c'est ce qui rend la
promesse énonçable en une phrase.

Le corollaire est que l'audience se choisit \textbf{hors trajet}, en modifiant la liste
Confiance --- un geste à froid, délibéré, et qui ne notifie personne, conformément à la
règle de confidentialité posée au volet~2. Au moment de partir, il n'y a plus rien à
décider~; c'est précisément l'objectif.
\end{note}

Une seule exception en cours de route, à l'échelle du trajet et sans effet sur la liste~: le
propriétaire peut \textbf{révoquer} un destinataire. C'est aussi le mécanisme qu'emprunte la
révocation automatique en cas de blocage.

\begin{alerte}[Deux défauts connus des listes deviennent des défauts de sûreté]
\code{deleteList} (\code{src/\allowbreak controllers/\allowbreak contactListController.js}) protège le \emph{nom} des
listes système mais pas leur \emph{existence}~: supprimer « Confiance » réussit, et
\code{ensureDefaultContactLists} la recrée \textbf{vide} au prochain
\mGET{}~\code{/api/contact-lists}. Aujourd'hui on perd un rangement~; demain on vide un
cercle de sécurité sans le dire.

\code{getListMembers} appelle \code{maskPresenceIfBlocked} par ligne, mais
\textbf{n'exclut pas} les membres qui ont bloqué le propriétaire. Le composeur de trajet
proposerait donc de confier sa sécurité à quelqu'un qui l'a bloqué.

Les deux figurent déjà aux points ouverts du volet~2. Ils deviennent des \textbf{prérequis
du même lot}.
\end{alerte}

\section{Mémoire et base}

\begin{note}[La règle de partage]
Un tic de position peut se perdre. Une transition d'état, non.
\end{note}

En mémoire, dans un \code{src/\allowbreak socket/\allowbreak state/\allowbreak tripState.js} calqué sur \code{callState.js}~:
dernière position, horodatage du dernier envoi, compteur de décimation, ensemble des
destinataires, ensemble des blocages, minuterie de péremption, cache du lien d'appareil.

En base, tout ce qui doit survivre à un redémarrage. C'est la différence majeure avec les
appels~: \textbf{un trajet survit au redémarrage du serveur}. La base fait foi~; la mémoire
n'est qu'un cache et un limiteur de débit, reconstruit paresseusement au premier contact.

\chapter{Le fil des positions}

\section{Trois régimes, pilotés par la distance}

Le moteur est \code{Geolocator.getPositionStream()} et son \code{distanceFilter}, pas une
minuterie. Le système d'exploitation pousse un point \emph{quand l'appareil a bougé}~; à
l'arrêt, il n'émet rien.

\begin{center}
\begin{tabularx}{\linewidth}{@{}p{2.3cm}p{3.6cm}p{1.7cm}p{1.2cm}p{1.5cm}X@{}}
\toprule
\thd{Régime} & \thd{Déclencheur} & \thd{Précision} & \thd{Filtre} & \thd{Plancher} &
\thd{Battement} \\
\midrule
Nominal  & à plus de 10 min de l'ETA         & \code{balanced} & 25 m &  5 s & 45 s \\
Approche & à moins de 10 min, ou 1 km du but & \code{high}     & 15 m &  4 s & 20 s \\
Alerte   & états \code{alert} et \code{sos}  & \code{best}     & 15 m &  5 s & 15 s \\
\bottomrule
\end{tabularx}
\end{center}

Le battement forcé est indispensable~: sans lui, on ne distingue pas \emph{immobile} de
\emph{traceur mort}. C'est exactement l'information dont un destinataire a besoin. Il ne
rallume pas le GPS~--- il renvoie le dernier point connu
(\code{Geolocator.\allowbreak getLastKnownPosition()}, déjà utilisé comme repli dans
\code{lib/\allowbreak screens/\allowbreak chats/\allowbreak location\_\allowbreak
picker\_\allowbreak screen.dart}) avec son horodatage de capture réel. L'écart entre
\code{recorded\_at} et \code{received\_at} dit alors, à lui seul, si l'appareil mesure
encore.

\begin{alerte}[Le \code{distanceFilter} seul ne borne pas le débit]
Un filtre de 50~m produit un point toutes les 36~s à pied, mais toutes les
\textbf{6 secondes} dans un taxi à 30~km/h, et toutes les 2~secondes sur voie rapide. Le
débit est donc proportionnel à la vitesse --- et le taxi, qui est le cas d'usage principal,
est le plus coûteux.

D'où le \textbf{plancher} de la quatrième colonne~: deux envois ne peuvent pas être séparés
de moins de $n$ secondes. Les points intermédiaires ne sont pas jetés, ils sont
\emph{groupés} dans l'envoi suivant --- la trace garde sa finesse, seul le nombre de paquets
est borné.
\end{alerte}

\section{Aucune de ces valeurs n'est écrite en dur}

Ce sont des \textbf{points de départ à régler en service}, pas des constantes. Deux
conséquences de conception, et la seconde est la moins évidente.

\textbf{Le serveur en est propriétaire.} Un module
\code{src/\allowbreak constants/\allowbreak tripPolicy.js} lit l'environnement et expose la
politique. Il n'existe pas de table de réglages globale dans Alanya
(\code{appSettingsService.js} est \emph{par utilisateur})~; la convention maison est la
variable d'environnement, comme \code{NOTIFICATION\_\allowbreak ANDROID\_\allowbreak
NATIVE\_V2}.

\begin{alerte}[La politique doit être \emph{servie}, jamais compilée dans l'application]
La cadence s'applique sur le téléphone. Si ces valeurs vivent dans des constantes Dart,
les changer impose une \textbf{publication de l'application} et des semaines d'adoption ---
autant dire qu'elles ne se règlent jamais.

\mPOST{}~\code{/api/trips} et \mGET{}~\code{/api/trips/active} renvoient donc un objet
\code{policy}, que le client applique. Les constantes Dart ne servent que de repli, si la
réponse est absente ou le serveur plus ancien. \code{trip:state} la reporte aussi, ce qui
permet à un réglage de prendre effet \textbf{sur un trajet déjà en cours}.
\end{alerte}

\begin{center}
\begin{tabularx}{\linewidth}{@{}p{5.4cm}p{1.5cm}X@{}}
\toprule
\thd{Variable} & \thd{Défaut} & \thd{Ce que coûte un réglage plus large} \\
\midrule
\code{TRIP\_FILTER\_M\_NOMINAL}    & 25 m  & Trace plus anguleuse en ville~; sans effet sur
                                            la sûreté. \\
\code{TRIP\_FLOOR\_S\_NOMINAL}     &  5 s  & Le pin saccade. C'est la molette la plus sûre à
                                            tourner~: elle ne change que le nombre de
                                            paquets. \\
\code{TRIP\_BEAT\_S\_NOMINAL}      & 45 s  & Un arrêt met plus longtemps à se confirmer, et
                                            le seuil de péremption recule d'autant. \\
\code{TRIP\_FILTER\_M\_APPROACH}   & 15 m  & L'arrivée se détecte plus tard. \\
\code{TRIP\_FLOOR\_S\_APPROACH}    & 4 s   & --- \\
\code{TRIP\_BEAT\_S\_APPROACH}     & 20 s  & --- \\
\code{TRIP\_FILTER\_M\_ALERT}      & 15 m  & \textbf{À ne pas relâcher} --- c'est le seul
                                            régime où la précision sert vraiment. \\
\code{TRIP\_FLOOR\_S\_ALERT}       & 5 s   & \textbf{À ne pas relâcher.} \\
\code{TRIP\_BEAT\_S\_ALERT}        & 15 s  & \textbf{À ne pas relâcher.} \\
\code{TRIP\_STALE\_FACTOR}         & 3     & Seuil de péremption $=$ facteur $\times$
                                            battement $+$ 30 s. \\
\code{TRIP\_BROADCAST\_MIN\_S}     & 5 s   & Cadence de rediffusion aux proches. \\
\code{TRIP\_PERSIST\_MIN\_S}       & 30 s  & Finesse de la trace rejouable, et volume de
                                            \code{trip\_point}. \\
\code{TRIP\_LOW\_BATTERY\_PCT}     & 15 \% & Seuil de bascule en suivi ralenti. \\
\bottomrule
\end{tabularx}
\end{center}

\begin{note}[Deux familles de réglages, qui ne se propagent pas pareil]
Les valeurs de \textbf{cadence} ci-dessus sont pures dépenses~: elles s'appliquent à chaud,
y compris sur un trajet en cours.

Les valeurs de \textbf{contrat} --- \code{dest\_radius\_m}, \code{grace\_minutes}, les
relances, la durée maximale --- sont d'une autre nature~: elles définissent ce qui a été
promis à l'utilisateur au moment du départ. Elles sont donc \textbf{recopiées dans la ligne
\code{trip}} à la création (les deux premières y sont déjà des colonnes), et un réglage
ultérieur ne modifie \textbf{jamais} un trajet déjà parti. Sinon on changerait le contrat
après signature.
\end{note}

\textbf{Garde batterie.} Sous 15~\%, retour au régime nominal avec un battement de 120~s, et
le cercle est \emph{informé} (« batterie faible --- suivi ralenti »). Sous 5~\%, émission
d'un événement portant la dernière position~: un téléphone qui meurt est un signal utile, et
il ne doit pas être confondu avec une disparition.

\section{Décimation à trois étages}

\begin{enumerate}[leftmargin=1.6em]
  \item \textbf{Client} --- le \code{distanceFilter} supprime déjà tout ce qui n'est pas un
        déplacement.
  \item \textbf{Socket} --- le serveur ne rediffuse pas plus d'une position toutes les 5~s
        par trajet, quelle que soit la cadence entrante.
  \item \textbf{Base} --- une insertion dans \code{trip\_point} toutes les 30~s au plus, par
        lots de 10~s. La dernière position vit dans les colonnes \code{last\_*} de
        \code{trip}, qui sont écrasées~; la table \code{trip\_point} n'est là que pour
        rejouer la trace.
\end{enumerate}

\section{Détection d'arrivée}

Trois conditions cumulatives, et une conséquence~:

\begin{enumerate}[leftmargin=1.6em]
  \item être à moins de \code{dest\_radius\_m} (100~m par défaut)~;
  \item y rester \textbf{60~s} --- l'hystérésis~;
  \item se déplacer à moins de $\sim$15 km/h --- le portillon de vitesse.
\end{enumerate}

Les relevés dont \code{acc\_m > 100} sont \textbf{ignorés pour la détection}, mais affichés
quand même, grisés, avec leur cercle de précision. Mieux vaut dire « à 60~m près » que
planter un point mensonger.

\begin{note}[Pourquoi 100 m, et pas 50]
En ville, la précision courante est de 20 à 65~m, et le multi-trajet en canyon urbain
dépasse régulièrement 100~m. Un rayon de 50~m ne déclencherait jamais en centre-ville~; 500~m
déclencherait à un pâté de maisons. Le défaut est donc \textbf{100~m} (compromis entre
détection fiable et question trop précoce)~; l'hystérésis et le portillon de vitesse font le
reste du travail~: passer devant sa rue sans s'arrêter est le faux positif numéro un.
\end{note}

Et la conséquence, déjà énoncée mais qui mérite d'être répétée ici~: \textbf{une arrivée
détectée ne clôt rien}. Elle fait passer en \code{awaiting\_confirm} et pose la question.

\section{Coût réseau et charge serveur}

Une position pèse environ 130 octets de JSON, et $\sim$220 octets sur le fil une fois
l'enveloppe Socket.IO, WebSocket, TLS et TCP/IP comptée. Ce n'est \textbf{pas} un appel
HTTP~: la trame part sur la socket déjà ouverte --- celle des messages, de la présence et
des appels --- sans en-tête, sans jeton à revalider, et \textbf{sans réponse attendue}.

C'est la capture qui suit la vitesse~; le plancher, lui, borne les envois~:

\begin{center}
\begin{tabularx}{\linewidth}{@{}p{3.4cm}p{2.6cm}p{2.4cm}X@{}}
\toprule
\thd{Allure} & \thd{Captures/h} & \thd{Envois/h} & \thd{Montant, trajet de 30 min} \\
\midrule
5 km/h, à pied      &   100 &   100 & $\sim$11 ko \\
30 km/h, taxi ville &   600 &   240 & $\sim$26 ko \\
90 km/h, voie rapide & 1\,800 &   240 & $\sim$26 ko \\
\bottomrule
\end{tabularx}
\end{center}

Le plancher de 15~s plafonne les envois à 240 par heure quelle que soit l'allure --- sans
lui, un taxi en ville en produirait 600, et le double sur voie rapide.

En descendant, un proche \emph{dont l'écran est ouvert} reçoit au plus une trame toutes les
5~s, soit $\sim$160 ko/h. En pratique il ne regarde presque jamais~: il reçoit la
notification et ouvre la carte une ou deux fois.

\begin{alerte}[Le vrai poste de dépense est la carte, pas les positions]
Une session de suivi de cinq minutes avec une carte qui se déplace charge 60 à 150 tuiles,
soit \textbf{1 à 3 Mo}. Les tuiles coûtent \textbf{20 à 50 fois} plus cher que le flux de
positions qu'elles servent à afficher.

C'est ce qui justifie deux choix du dossier~: la carte du proche est une \emph{vignette}
figée dans la conversation, jamais une carte vivante~; et un fournisseur de tuiles sous
contrat est un prérequis de mise en service.
\end{alerte}

\begin{center}
\begin{tabularx}{\linewidth}{@{}p{5.6cm}X@{}}
\toprule
\thd{À 1 000 trajets simultanés} & \thd{Ordre de grandeur} \\
\midrule
Trames socket entrantes      & $\sim$70 / s --- aucune connexion nouvelle, elles empruntent
                               les sockets existantes \\
Émissions sortantes          & $\sim$100 à 200 / s \\
Insertions \code{trip\_point} & 33 lignes/s, groupées en un \code{INSERT} multi-lignes
                               toutes les 2~s \\
Stockage                     & 1 000 trajets de 30 min $=$ 60 000 lignes $\approx$ 6 Mo,
                               purgés à 24~h \\
Lignes de \code{job\_queue}   & 4 000 armées \\
\bottomrule
\end{tabularx}
\end{center}

Pour situer cet ordre de grandeur dans l'existant~: \code{typing:start} est aujourd'hui émis
\textbf{à chaque caractère frappé}, et coûte « 1+2N requêtes SQL par frappe »
(\code{docs/\allowbreak AUDIT\_\allowbreak SCALABILITE\_\allowbreak 2026-08-06.md}). Une
seule conversation de groupe animée coûte plus de SQL que cent trajets.

\begin{alerte}[Le plafond réel de la file est de 20 jobs par seconde]
\code{startJobWorker()} tourne à \code{setInterval(tick, 1000)} avec
\code{BATCH\_SIZE = 20} (\code{src/\allowbreak services/\allowbreak jobQueue.js}). Aux heures
où les ETA s'agglutinent --- 18\,h--20\,h --- c'est juste. \code{idx\_job\_pret} et
\code{FOR UPDATE SKIP LOCKED} tiennent la concurrence~; c'est \code{BATCH\_SIZE} qui est la
molette. À mesurer avant la montée en charge, pas après.
\end{alerte}

\begin{note}[Batterie --- le poste qui compte vraiment]
\code{balanced} coûte 3 à 6~\%/h, \code{high} 8 à 15~\%/h, \code{best} en continu 15 à
25~\%/h. Un taxi de 30 minutes en régime nominal revient à 2 ou 3~\%. C'est le régime
d'alerte qui est cher --- et c'est une des raisons du plafond de durée.
\end{note}


\chapter{L'échéance et l'escalade}

\section{L'escalade en trois temps}

\begin{center}
\begin{tabularx}{\linewidth}{@{}p{2.0cm}p{4.4cm}X@{}}
\toprule
\thd{Moment} & \thd{Job} & \thd{Effet} \\
\midrule
$T-5$ min & \code{trip\_eta\_soon}     & Silencieux, \textbf{propriétaire seul}~:
                                         « vous arrivez bientôt~? Prolonger / Je suis
                                         arrivé ». \\
$T$       & \code{trip\_eta\_due}      & État \code{awaiting\_confirm}, notification
                                         prioritaire \textbf{au propriétaire seul}~;
                                         relances à $+0$, $+3$ et $+7$ min. \\
$T+10$ min & \code{trip\_grace\_expired} & État \code{alert}~: le cercle est prévenu, avec
                                          la dernière position connue. \\
\bottomrule
\end{tabularx}
\end{center}

\begin{constat}[Les cinq types de notification poussée]
Chaque étape de l'escalade porte son propre type, et chacun compose son \emph{titre} et
son \emph{corps} côté serveur~: les envois sont \emph{data-only}, et quand l'application
est tuée c'est le système qui affiche~--- il n'y a personne pour traduire. Une charge
utile sans ces deux champs produit une notification parfaitement silencieuse.

\code{trip\_eta\_soon}, \code{trip\_due} et \code{trip\_reminder} vont au
\textbf{propriétaire seul}, sur le canal \code{alanya\_trip}, en priorité haute mais
\textbf{sans} traverser « Ne pas déranger »~: ce sont des rappels, pas des alertes.

\code{trip\_alert} et \code{trip\_sos} vont au cercle, sur le canal
\code{alanya\_trip\_alert}, dont l'usage audio est déclaré \code{alarm}. C'est ce qui
leur fait passer la catégorie « alarmes » de « Ne pas déranger », autorisée par défaut~---
la seule voie sans \code{ACCESS\_NOTIFICATION\_POLICY}, permission que l'utilisateur
doit accorder à la main. Sur iOS, \code{InterruptionLevel.timeSensitive} joue le même
rôle face au mode Concentration. Ce sont les deux seuls types de toute l'application
autorisés à traverser le silence.

Le \emph{tag} de notification dérive du trajet et non du type~: un trajet occupe une
ligne dans le tiroir, et le rappel d'échéance remplace le pré-avis au lieu de s'empiler à
côté. Même principe que la carte de conversation, réécrite sur place.
\end{constat}

\begin{note}[Pourquoi dix minutes, et non cinq]
En deçà de cinq minutes, chaque embouteillage devient une alerte, et le cercle apprend à
ignorer les alertes. C'est le véritable mode de défaillance de ce genre de produit~: pas
l'alerte manquée, mais l'alerte qu'on n'ouvre plus. Au-delà d'un quart d'heure, le délai
n'a plus d'utilité dans un incident réel.
\end{note}

La \textbf{prolongation} est à deux appuis, $+15$ ou $+30$ minutes, sans limite de nombre.
Chaque prolongation ré-arme les jobs et écrit un \code{trip\_event} visible du cercle
(« Awa a prolongé de 15 minutes »)~: cela rassure sans alerter.

\section{Les huit \code{kind} de jobs}

\code{trip\_eta\_soon}, \code{trip\_eta\_due}, \code{trip\_reminder},
\code{trip\_grace\_expired}, \code{trip\_max\_duration}, \code{trip\_no\_watcher},
\code{trip\_stale\_check}, \code{trip\_points\_purge}. Un service
\code{src/\allowbreak services/\allowbreak tripWorkers.js} les enregistre, sur le modèle de
\code{src/\allowbreak services/\allowbreak broadcastWorkers.js}.

\code{dedupe\_key = "trip\_<id>"} pour chacun. La contrainte
\code{UNIQUE uq\_\allowbreak job\_\allowbreak dedupe (kind, dedupe\_\allowbreak key)} de \code{migrations/\allowbreak 039\_job\_queue.sql}
rend le ré-armement trivial~: une transaction qui supprime les lignes du trajet puis
appelle \code{enqueue} avec le nouveau \code{run\_after}.

\begin{alerte}[\code{enqueue} est silencieux sur collision]
\code{enqueue} fait \code{ON DUPLICATE KEY UPDATE id = id} et \textbf{renvoie \code{null}}
si la clé de déduplication existe déjà. Prolonger un trajet sans supprimer d'abord la ligne
existante \emph{paraît} fonctionner et laisse l'ancienne échéance armée --- l'alerte
partirait à l'heure initiale. Le \code{DELETE} préalable n'est pas une optimisation, c'est
la correction.
\end{alerte}

\begin{note}[La file est \emph{at-least-once}]
Un job peut se déclencher après la clôture du trajet --- reprise après incident, orphelin
récupéré par \code{reclaimOrphans} au bout de dix minutes. \textbf{Chaque handler relit
\code{trip.state} et sort en silence si le trajet est terminal.} C'est la seule protection
qui tienne.
\end{note}

La purge nocturne tourne sous le bail \code{trip\_nightly\_purge} via
\code{withLease()} (\code{src/\allowbreak services/\allowbreak schedulerLease.js}), à côté de
\code{broadcast\_nightly\_purge} et \code{welcome\_status\_purge}.

\chapter{API}

\section{REST --- \code{src/routes/trips.js}}

Style copié de \code{contactListController.js}~: cadrage sur \code{req.user.alanyaID}, jamais
sur un identifiant fourni par le client~; \code{pool.execute}~; codes 400/403/404/409.

\begin{center}
\begin{tabularx}{\linewidth}{@{}p{5.6cm}X@{}}
\toprule
\mPOST{}~\code{/api/trips}                    & Démarrer. \code{\{clientId, kind, dest?,
                                                etaAt? | durationMin?, note?\}} ---
                                                l'audience n'est pas un paramètre \\
\mPOST{}~\code{/api/trips/sos}                & SOS autonome --- crée le trajet \\
\mGET{}~\code{/api/trips/active}              & Reprise à froid~: mon trajet, et ceux que je
                                                suis \\
\mGET{}~\code{/api/trips/:id}                 & Détail et frise d'événements \\
\mGET{}~\code{/\allowbreak api/\allowbreak trips/\allowbreak :id/\allowbreak points?since=}   & La trace --- vide après purge \\
\mPOST{}~\code{/api/trips/:id/extend}         & \code{\{minutes\}}, ré-arme les jobs \\
\mPOST{}~\code{/api/trips/:id/confirm}        & Arrivée confirmée \\
\mPOST{}~\code{/api/trips/:id/cancel}         & Arrêt avant échéance \\
\mPOST{}~\code{/api/trips/:id/sos}            & Escalade d'un trajet existant \\
\mDEL{}~\code{/\allowbreak api/\allowbreak trips/\allowbreak :id/\allowbreak watchers/\allowbreak :id}    & Révoquer un destinataire \\
\mPOST{}~\code{/api/trips/:id/seen}           & Acquittement d'un destinataire \\
\mGET{}~\code{/\allowbreak api/\allowbreak trips/\allowbreak history?cursor=}     & Historique du propriétaire \\
\mDEL{}~\code{/api/trips/:id}                 & Suppression par le propriétaire \\
\bottomrule
\end{tabularx}
\end{center}

Codes d'erreur~: \code{TRUST\_LIST\_EMPTY} 409, \code{TRIP\_ALREADY\_ACTIVE} 409,
\code{TRIP\_TERMINAL} 409, \code{TRIP\_INCIDENT\_LOCKED} 423,
\code{DEVICE\_NOT\_OWNER} 403.

L'audience n'étant pas transmise par le client, elle n'a pas de code d'erreur propre~: le
serveur lit la liste \code{trust} du porteur du jeton, et refuse en 409
\code{TRUST\_LIST\_EMPTY} si elle est vide.

\begin{note}[Toute réponse portant un trajet porte aussi la politique]
\mPOST{}~\code{/api/trips}, \mGET{}~\code{/api/trips/active} et l'événement
\code{trip:state} renvoient un objet \code{policy}~: les trois régimes, avec leurs filtres,
planchers et battements, tels que \code{tripPolicy.js} les expose au moment de l'appel. Le
client l'applique~; ses constantes Dart ne servent que si le champ manque.

C'est la seule façon de régler la cadence sans publier l'application. Le champ doit donc
exister \textbf{dès la première version}, même s'il ne renvoie d'abord que les valeurs par
défaut --- l'ajouter plus tard imposerait exactement la publication qu'on cherche à éviter.
\end{note}

\begin{note}[Aucune position ne passe par REST en régime normal]
Un aller-retour HTTP tous les 50 mètres serait absurde. Les positions montent par socket~;
\mGET{}~\code{/api/trips/:id/points} ne sert qu'à rejouer une trace, et le lot de secours
n'existe que pour vider le tampon hors ligne.
\end{note}

\section{Socket.IO --- \code{src/\allowbreak socket/\allowbreak handlers/\allowbreak trips.js}}

Enregistré à côté de \code{calls.js} et \code{meetings.js} dans \code{server.js}. Room
\code{trip\_<id>}, rejointe après vérification d'appartenance, exactement comme
\code{join\_conversation} vérifie \code{conv\_participants}.

\begin{center}
\begin{tabularx}{\linewidth}{@{}p{5.0cm}X@{}}
\toprule
\thd{Client $\rightarrow$ serveur} & \thd{Objet} \\
\midrule
\code{trip:subscribe}      & Autorise, puis renvoie l'état complet \\
\code{trip:unsubscribe}    & --- \\
\code{trip:position}       & Un point. Rejeté si \code{deviceId} n'est pas
                             \code{trip.owner\_device} \\
\code{trip:position\_batch} & Vidange du tampon hors ligne, 200 points au plus \\
\code{trip:claim\_device}  & Reprendre la main depuis un autre appareil \\
\code{trip:signal}         & \code{gps\_off}, \code{permission\_denied},
                             \code{background\_killed}, \code{battery\_saver} \\
\code{trip:seen}           & Acquittement d'un destinataire \\
\bottomrule
\end{tabularx}
\end{center}

\begin{center}
\begin{tabularx}{\linewidth}{@{}p{5.0cm}X@{}}
\toprule
\thd{Serveur $\rightarrow$ client} & \thd{Objet} \\
\midrule
\code{trip:state}          & État complet --- à la souscription et à chaque transition \\
\code{trip:position}       & Fan-out room, plafonné à une émission / 5 s / trajet \\
\code{trip:card\_update}   & Mutation de la carte de type 9 \\
\code{trip:alert}          & Le trajet et sa dernière position \\
\code{trip:stale}          & Plus de position reçue \\
\code{trip:signal}         & Relais du motif déclaré par le propriétaire \\
\code{trip:device\_revoked} & Un autre appareil a repris la main \\
\code{trip:closed}         & Fin, avec son motif \\
\bottomrule
\end{tabularx}
\end{center}

\begin{alerte}[La room ne porte que le mouvement --- jamais l'alerte]
\code{trip\_<id>} sert \textbf{uniquement} au flux haute fréquence. Tout ce qui doit
\emph{atteindre} un destinataire --- état, alerte, clôture --- passe par
\code{emitToUser} (\code{src/\allowbreak utils/\allowbreak userSocketRegistry.js}) \textbf{plus} la notification
poussée. Un destinataire dont aucun appareil n'est abonné à la room doit quand même recevoir
l'alerte. \textbf{Une alerte ne transite jamais par la seule room.}
\end{alerte}

\code{isUserOnline()} et \code{hasForegroundSocket()} décident si un destinataire reçoit le
fan-out socket ou seulement la notification~: on ne diffuse pas dans une room que personne ne
regarde.

\section{Le verrou d'appareil}

\code{trip.owner\_device} porte l'identifiant de l'unique appareil autorisé à émettre. Les
autres appareils du compte affichent « trajet en cours sur votre autre appareil » et peuvent
reprendre la main explicitement~: \code{trip:claim\_device} bascule la colonne, l'ancien
porteur reçoit \code{trip:device\_revoked} et arrête son service.

C'est le motif de \code{src/\allowbreak socket/\allowbreak state/\allowbreak callDeviceOwnership.js}, transposé --- avec la même
justification qu'aux appels~: sans lui, deux appareils entrelacent leurs positions et la
trace zigzague entre deux endroits.

\chapter{Le message de type 9}

\section{Pourquoi un type de message}

Les types existants vont de 0 à 8~; le 9 est libre. Ils ne sont d'ailleurs documentés nulle
part au complet~: le commentaire de \code{LocalMessages.type}
(\code{lib/\allowbreak core/\allowbreak db/\allowbreak app\_database.dart}) s'arrête à \code{5=localisation}, le 6 et le 8 sont
des constantes nommées (\code{kSystemMessageType} dans
\code{lib/\allowbreak core/\allowbreak utils/\allowbreak system\_\allowbreak event\_\allowbreak payload.dart}, \code{kWelcomeCtaMessageType} dans
\code{welcome\_cta\_payload.dart}), et le 7 --- la carte de contact --- n'existe qu'en
littéral, dans \code{chat\_bubbles.dart} et \code{talky\_models.dart}. Le type 9 mérite une
constante nommée et un commentaire tenu à jour.

Le choisir plutôt que d'inventer une entité de
conversation apporte trois choses gratuitement~: la notification poussée, le compteur de
non-lus, et la persistance dans l'archive. Et \code{local\_messages} n'a pas à changer --- le
JSON tient dans \code{content}, comme pour le type 5.

\section{L'état dans le message, le mouvement dans le socket}

C'est le point technique central du volet.

\begin{note}[Cinq écritures, pas cinq mille]
Le \code{content} du message porte l'\textbf{état}~: il est réécrit uniquement aux
transitions --- cinq à six fois sur toute la vie d'un trajet. Les \emph{positions} ne
touchent jamais la table \code{message}. Entre deux transitions, la carte s'anime à partir
des \code{trip:position} reçus.
\end{note}

La réécriture passe par \code{trip:card\_update}, et \textbf{jamais} par
\code{message:send}~: ce dernier remonterait \code{sentAt}, incrémenterait
\code{unreadCount} et relancerait une notification à chaque changement d'état.

\textbf{Exception --- les incidents.} Une alerte ou un SOS écrit \emph{en plus} un message
système de type 6, pour que l'incident reste dans l'archive même si la carte est refermée ou
le trajet purgé. Deux lignes au maximum par trajet, dans le pire cas.

\section{Compatibilité descendante}

\begin{alerte}[Un client ancien afficherait « Média »]
Les replis existent des deux côtés --- \code{default:} dans \code{\_buildMedia}
(\code{lib/\allowbreak screens/\allowbreak chats/\allowbreak chat/\allowbreak chat\_bubbles.dart}) et
\code{default: return mediaName ?? 'Média'} dans \code{messagePreview}
(\code{src/\allowbreak utils/\allowbreak messagePreview.js}) --- mais ils rendent un libellé générique et, dans le
pire cas, le JSON brut. Le rendu dégradé (« Trajet de confiance --- mettez à jour
l'application ») et l'aperçu de conversation doivent être écrits dans les deux fichiers
\textbf{avant} la première émission d'un type 9 en production.
\end{alerte}

L'aperçu suit le patron de \code{locationPreviewFromContent}~: un court-circuit
\code{if (t === 9)} dans \code{messagePreview}, et une entrée dans \code{mediaTypeLabel}. Le
commentaire existant du fichier vaut aussi ici --- \emph{ne jamais exposer le
\code{content} brut}.

\chapter{Notifications}

Nouveaux types dans \code{src/\allowbreak notifications/\allowbreak notificationContract.js}~:
\code{trip\_started}, \code{trip\_alert}, \code{trip\_sos}, \code{trip\_closed},
\code{trip\_watcher\_seen}, et un constructeur \code{buildTripPayload} en
\code{schemaVersion: '2'} portant \code{tripId}, \code{state}, \code{ownerName},
\code{lastLat} / \code{lastLng} / \code{lastAt} et un lien profond.

\begin{center}
\begin{tabularx}{\linewidth}{@{}p{3.4cm}p{2.2cm}p{1.6cm}X@{}}
\toprule
\thd{Type} & \thd{Priorité} & \thd{TTL} & \thd{Rendu} \\
\midrule
\code{trip\_started}      & normale & 1 h   & Notification ordinaire \\
\code{trip\_alert}        & haute   & 120 s & \emph{data-only}, plein écran \\
\code{trip\_sos}          & haute   & 120 s & \emph{data-only}, plein écran, son propre \\
\code{trip\_closed}       & basse   & 1 h   & Silencieuse \\
\code{trip\_watcher\_seen} & basse  & 1 h   & Silencieuse, propriétaire seul \\
\bottomrule
\end{tabularx}
\end{center}

Le TTL court des alertes est délibéré~: une alerte de sûreté vieille de trois heures est du
bruit, et son affichage tardif décrédibilise les suivantes. Le modèle est
\code{sendCallToUser} / \code{buildIncomingCallPayload}, déjà en \code{ttl: 45\_000}.

\begin{alerte}[Contournement de « Ne pas déranger » --- pour deux types seulement]
\code{trip\_alert} et \code{trip\_sos} passent outre les préférences de notification et le
mode « Ne pas déranger ». Une alerte de sûreté qu'un réglage de silence peut étouffer n'est
pas une alerte de sûreté. Le contournement s'arrête là~: \code{trip\_started} est une
notification ordinaire et \code{trip\_closed} est silencieuse. À valider explicitement --- il
engage le produit.
\end{alerte}

Côté client, un canal Android dédié \code{alanya\_trip\_alert} en importance haute, avec son
propre son et un \code{fullScreenIntent} pour le SOS~; le canal d'appel
(\code{alanya\_call\_media}) donne le patron.

\chapter{Client Flutter}

\section{Cache local --- Drift 23 $\rightarrow$ 24}

\code{schemaVersion} vaut aujourd'hui \textbf{23} (\code{lib/\allowbreak core/\allowbreak db/\allowbreak app\_database.dart}),
les listes de contacts l'ayant porté là en trois étapes. Les trajets prennent le \textbf{24}~:

\begin{lstlisting}[language=Dart, numbers=none]
if (from < 24) {
  await m.createTable(localTrips);
  await m.createTable(localTripPoints);
  await m.createTable(localTripEvents);
}
\end{lstlisting}

Pure création de tables, aucune migration de données. \code{local\_messages} ne change pas.
Puis \code{dart run build\_runner build -{}-delete-conflicting-outputs}. La montée doit
passer par \code{onUpgrade}, jamais par une recréation.

\code{LocalTripPoints} est un \textbf{anneau borné}~: tampon hors ligne côté propriétaire,
cache de polyligne côté destinataire. Le plafonnement (suppression des plus anciens au-delà
de $N$ par trajet) vit dans le DAO, pas dans le schéma.

Deux nouveaux services~: \code{lib/\allowbreak core/\allowbreak services/\allowbreak trip\_repository.dart} (REST et Drift) et
\code{lib/\allowbreak core/\allowbreak services/\allowbreak trip\_\allowbreak socket\_\allowbreak service.dart}. Les deux tables sont à ajouter à
\code{clearSession()}.

\section{Le service de localisation en avant-plan (Android)}

\code{TripLocation\allowbreak ForegroundService.kt}, cloné de \code{CallMedia\allowbreak ForegroundService.kt}~:
canal \code{alanya\_trip\_location}, identifiant de notification distinct, et le même
conditionnement de \code{foregroundServiceType} au SDK 34 que l'existant. Pont
\code{TripLocationBridge.kt} sur le canal \code{com.alanya237.alanya/\allowbreak trip\_location},
calqué sur \code{CallMediaBridge.kt}.

\begin{note}[\code{START\_STICKY}, et non \code{START\_NOT\_STICKY}]
Le service d'appel est délibérément non collant~: un appel tué ne doit pas ressusciter. Un
trajet, si. Le service relit le trajet ouvert dans Drift au redémarrage et reprend l'émission.
\end{note}

Orchestrateur Dart \code{lib/\allowbreak core/\allowbreak services/\allowbreak trip\_session\_guard.dart}, cloné de
\code{lib/\allowbreak core/\allowbreak services/\allowbreak call\_session\_guard.dart}~: singleton, compteur de références,
\code{WidgetsBindingObserver}, démarrage et arrêt du service, propriété de l'abonnement
\code{getPositionStream()}, application des trois régimes, gestion du tampon hors ligne.

\begin{alerte}[Le contenu de la notification persistante a une portée éthique]
Elle doit \textbf{nommer les destinataires} (« Partagé avec Maman, Karim ») et offrir un
« Arrêter » en un appui. Une notification de suivi qui ne dit pas à qui elle envoie la
position, c'est le comportement d'un logiciel espion --- et c'est ce qui distingue cette
fonctionnalité d'un traceur.
\end{alerte}

\section{Permissions à ajouter}

\begin{center}
\begin{tabularx}{\linewidth}{@{}p{2.0cm}p{5.6cm}X@{}}
\toprule
\thd{Cible} & \thd{Clé} & \thd{État actuel} \\
\midrule
Android & \code{ACCESS\_\allowbreak BACKGROUND\_\allowbreak LOCATION}   & absente \\
Android & \code{FOREGROUND\_\allowbreak SERVICE\_\allowbreak LOCATION}  & absente \\
Android & \code{<service … foregroundServiceType="location">} & absent \\
iOS & \code{NSLocation\allowbreak Always\allowbreak AndWhenInUse\allowbreak UsageDescription} & absente \\
iOS & \code{UIBackgroundModes: location}        & absente --- la liste vaut
                                                  \code{audio}, \code{remote-notification},
                                                  \code{voip} \\
\bottomrule
\end{tabularx}
\end{center}

Sont déjà là et réutilisables~: \code{ACCESS\_FINE\_LOCATION},
\code{ACCESS\_COARSE\_LOCATION}, \code{FOREGROUND\_SERVICE}, \code{WAKE\_LOCK},
\code{POST\_NOTIFICATIONS}, \code{REQUEST\_\allowbreak IGNORE\_\allowbreak BATTERY\_\allowbreak OPTIMIZATIONS}.

\section{iOS --- ce qui est garanti, et ce qui ne l'est pas}

\code{allowsBackgroundLocationUpdates = true},
\code{showsBackgroundLocationIndicator = true} --- \textbf{on garde la pastille bleue}, on ne
la masque jamais --- et \code{pausesLocationUpdatesAutomatically = false}. Repli sur
l'enregistrement des changements significatifs de position.

\begin{note}[Ne pas promettre la parité]
La garantie iOS est « au mieux, avec péremption explicite ». Ce n'est pas un aveu de
faiblesse~: c'est la conséquence directe du postulat du chapitre~1. Le filet est la chaîne
d'échéance côté serveur, pas le flux de positions --- et cette chaîne est identique sur les
deux plateformes.
\end{note}

\section{Interface}

\begin{center}
\begin{tabularx}{\linewidth}{@{}p{6.2cm}X@{}}
\toprule
\thd{Fichier} & \thd{Objet} \\
\midrule
\code{screens/\allowbreak trips/\allowbreak trips\_hub\_screen.dart} \emph{(nouveau)} &
  Démarrer, trajet en cours, historique. \\
\code{screens/\allowbreak trips/\allowbreak trip\_\allowbreak compose\_\allowbreak screen.dart} \emph{(nouveau)} &
  Type, destination, heure, destinataires, consentement. \\
\code{screens/\allowbreak trips/\allowbreak trip\_live\_screen.dart} \emph{(nouveau)} &
  Suivi --- vue propriétaire et vue destinataire. \\
\code{widgets/\allowbreak chat/\allowbreak trip\_message\_card.dart} \emph{(nouveau)} &
  La carte de type 9, huit états. \\
\code{widgets/\allowbreak trips/\allowbreak trip\_banner.dart} \emph{(nouveau)} &
  Bandeau persistant au-dessus de la barre de navigation. \\
\code{screens/\allowbreak profile/\allowbreak profile\_screen.dart} &
  Entrée « Trajets de confiance », à côté des listes. \\
\code{screens/\allowbreak chats/\allowbreak chat/\allowbreak chat\_actions.dart} &
  Le menu position devient un choix à deux entrées. \\
\code{lib/l10n/*} & Libellés des états, de l'escalade et du consentement. \\
\bottomrule
\end{tabularx}
\end{center}

\begin{constat}[Aucun sixième onglet --- et pas besoin]
\code{lib/\allowbreak screens/\allowbreak home/\allowbreak glass\_nav\_bar.dart} construit sa rangée par
\code{List.generate(5, …)}~: Discussions, Appels, Statuts, Réunions, Profil. Le bandeau
persistant se pose \emph{au-dessus} de la barre, dans l'espace déjà réservé par
\code{kGlassNavBarSpace}. Il est visible sur les cinq onglets~: c'est le sixième onglet sans
en être un.
\end{constat}

\begin{alerte}[Les tuiles cartographiques ne sont pas gratuites à cette échelle]
\code{location\_\allowbreak picker\_\allowbreak screen.dart} charge \code{tile.openstreetmap.org} une fois, pour un
écran. Un trajet de trente minutes suivi par cinq personnes en charge des \emph{centaines}.
La politique d'usage d'OpenStreetMap l'interdit explicitement pour ce volume. Un fournisseur
sous contrat est un \textbf{prérequis de mise en production}, pas une optimisation ultérieure.
Le composant \code{flutter\_map} ne change pas --- seule l'URL des tuiles change.
\end{alerte}

\chapter{Vie privée, coercition et administration}

\section{Base légale et minimisation}

La base légale est le \textbf{consentement}, donné par trajet, révocable à tout instant en
arrêtant le partage. Pas l'intérêt légitime~: une trace de déplacement continue liée à une
identité révèle par inférence la pratique religieuse, l'état de santé et l'orientation
sexuelle. Elle n'est pas une donnée personnelle ordinaire.

Le géocodage inverse passe par \code{nominatim.openstreetmap.org}
(\code{lib/\allowbreak core/\allowbreak utils/\allowbreak location\_payload.dart}) --- un tiers. On géocode donc
\textbf{la destination une seule fois}, à la création, et \textbf{jamais la trace}.

\section{Rétention --- le tableau fait foi}

\begin{center}
\begin{tabularx}{\linewidth}{@{}p{3.4cm}p{3.2cm}p{3.2cm}X@{}}
\toprule
\thd{Donnée} & \thd{Trajet normal} & \thd{Trajet à incident} & \thd{Raison} \\
\midrule
\code{trip\_point} & 24 h après clôture & 30 jours & Un historique permanent des
déplacements de tous les utilisateurs n'est justifié par aucun besoin produit. En cas
d'incident, la trace peut être une preuve. \\
\code{trip}        & 12 mois & 12 mois & L'historique produit se contente du fait. \\
\code{trip\_event} & 12 mois & 12 mois & Source de la frise. \\
\bottomrule
\end{tabularx}
\end{center}

La suppression par l'utilisateur est immédiate pour un trajet normal, avec effacement de la
carte chez les destinataires.

\begin{note}[Le verrou de 30 jours sur les incidents est une mesure anti-coercition]
Un trajet clos en \code{alert} ou \code{sos} ne peut pas être supprimé avant 30 jours
(423 \code{TRIP\_INCIDENT\_LOCKED}). C'est délibéré~: un agresseur ne doit pas pouvoir
contraindre quelqu'un à effacer la trace. À expliquer dans l'interface, pas à subir.
\end{note}

La suppression de compte doit être complétée dans
\code{src/\allowbreak services/\allowbreak accountDeletionService.js}~: les \code{ON DELETE CASCADE} couvrent
l'essentiel, la vérification explicite reste due.

\section{Ce que l'administration voit --- et ce qu'on refuse de construire}

L'administration reçoit une page de \textbf{compteurs agrégés} et rien d'autre~: trajets
démarrés, clos, en alerte par jour~; durée médiane~; taux d'alerte~; taux de SOS. Sans
\code{alanyaID}, sans coordonnées, sans libellé de destination, sans identité de
destinataire. La coque \code{components/geolocation/} de l'administration Next.js
(\code{WorldMap.tsx}, \code{GeoStats.tsx}) se réutilise telle quelle pour cette vue.

\begin{note}[Deux options examinées, puis écartées]
Un \textbf{registre d'incidents pseudonymisé} (qui, quand, quelle issue) aiderait le
support~; un \textbf{bris de glace} sur la dernière position, à quatre yeux, avec
justification écrite, expiration à 15 minutes, journal en ajout seul et notification à la
personne concernée, aiderait en cas d'incident réel.

Les deux sont écartés pour cette étape. Non parce qu'ils seraient mal conçus, mais parce que
la posture la plus défendable est celle où \textbf{la donnée n'est pas accessible}. Et la
rétention l'emporte de toute façon sur le contrôle d'accès~: à 24~h de conservation hors
incident, la fenêtre de consultation n'existe quasiment pas.

Si le besoin revient, la forme est connue --- une table \code{trip\_admin\_access} en ajout
seul, portant \code{admin\_id}, \code{approver\_id}, \code{reason}, \code{expires\_at} et
\code{subject\_notified\_at}. Elle n'est pas créée par \code{051\_trips.sql}.
\end{note}

\begin{alerte}[L'AIPD est un prérequis, pas un rattrapage]
Un suivi de localisation systématique relève de l'article 35 du RGPD~: l'analyse d'impact
est \textbf{obligatoire}. Elle doit être menée avant la mise en production, et elle doit
inclure un modèle de menace « violences conjugales » --- pas seulement un modèle de menace
technique.
\end{alerte}

\section{Détournement et coercition}

C'est le risque qui peut tuer la fonctionnalité, et il ne se traite pas par une case à
cocher. Sept garde-fous, traités comme des \emph{contraintes de conception}~:

\begin{enumerate}[leftmargin=1.6em]
  \item \textbf{Le propriétaire est toujours l'émetteur.} Aucune demande de position, aucune
        invitation à partager, aucun flux d'acceptation qu'un partenaire puisse transformer
        en permanence.
  \item \textbf{Tout trajet est fini et démarré à la main.} Pas de « partager 24 h », pas de
        récurrence, pas d'interrupteur permanent.
  \item \textbf{La notification persistante nomme l'audience} et ne peut pas être masquée
        (service en avant-plan Android, pastille bleue iOS).
  \item \textbf{Les destinataires n'ont pas d'historique.} On ne peut pas ouvrir « tous les
        trajets passés d'Awa ». Cela tue l'interrogatoire « où étais-tu mardi~?~».
  \item \textbf{L'audience se choisit à froid, et en silence.} On ne compose pas son cercle
        au moment de partir~: on l'a fait avant, dans ses listes de contacts, et cette
        modification ne notifie personne --- ni celui qui entre, ni celui qui sort. Sans
        cela, « tu m'as retiré » deviendrait un levier de pression exerçable en direct.
  \item \textbf{La sortie est sûre.} « Arrêter le partage » est toujours à un appui, et un
        arrêt avant échéance n'envoie qu'un message neutre. Si arrêter était visiblement
        suspect, arrêter deviendrait punissable.
  \item \textbf{Le verrou de suppression de 30 jours} sur les trajets clos en incident.
\end{enumerate}

\begin{alerte}[Ce sont des atténuations, pas une solution]
Aucune de ces mesures n'empêche quelqu'un de contraindre un proche à démarrer un trajet.
Elles rendent la contrainte moins rentable et moins durable. Il faut l'écrire tel quel dans
l'AIPD~: prétendre le contraire serait la plus dangereuse des erreurs de conception.
\end{alerte}

\section{Le blocage, et le cercle qui se vide}

Le blocage est la primitive de sûreté de l'application~; le respecter n'est pas négociable.
Un destinataire qui bloque le propriétaire passe en \code{state = 'revoked'} avec
\code{revoke\_\allowbreak reason = 'blocked'}, quitte la room, voit sa carte effacée, et
\textbf{ne reçoit plus rien --- alertes comprises}. Le propriétaire lit « un membre n'est
plus destinataire », jamais « X vous a bloqué »~: l'application ne révèle pas un blocage.

Le contrôle tourne \textbf{à la création} (filtre de l'instantané) \textbf{et à chaque
diffusion persistée} (alerte, clôture) --- pas seulement à la création, car le blocage peut
survenir en route et la room ne le sait pas. La requête symétrique existe déjà~:
\code{isBlockedEitherWay} (\code{src/utils/blockUtils.js}), et sa forme groupée dans
\code{src/\allowbreak utils/\allowbreak conversationParticipantsBatch.js}.

\begin{alerte}[S'il ne reste aucun destinataire, le trajet ne continue pas]
Un trajet sans destinataire n'est plus une fonctionnalité de sécurité~: c'est un traceur
personnel, un autre produit, et une responsabilité RGPD qu'on ne veut pas assumer. Le
propriétaire reçoit une invite immédiate et sonore~; sans réponse au bout de 5 minutes, le
job \code{trip\_no\_watcher} clôt en \code{closed\_unwatched} et la purge normale s'applique.
\end{alerte}

\chapter{Impact fichier par fichier}

\section{Backend}

\begin{center}
\begin{tabularx}{\linewidth}{@{}p{6.4cm}p{1.6cm}X@{}}
\toprule
\thd{Fichier} & \thd{Nature} & \thd{Objet} \\
\midrule
\code{migrations/051\_trips.sql}            & nouveau & Les quatre tables et le bail \\
\code{src/routes/trips.js}                  & nouveau & Les quatorze routes, Swagger inline \\
\code{src/\allowbreak controllers/\allowbreak tripController.js}    & nouveau & Machine à états, autorisations \\
\code{src/\allowbreak services/\allowbreak tripWorkers.js}          & nouveau & Les huit \code{kind} de jobs \\
\code{src/\allowbreak constants/\allowbreak tripPolicy.js}          & nouveau & Rayon, grâce, plafonds, cadences \\
\code{src/\allowbreak socket/\allowbreak handlers/\allowbreak trips.js}         & nouveau & Room, positions, verrou d'appareil \\
\code{src/\allowbreak socket/\allowbreak state/\allowbreak tripState.js}        & nouveau & Cache et limiteur de débit \\
\code{server.js}                            & modifié & Montage des routes et du handler \\
\code{src/\allowbreak utils/\allowbreak messagePreview.js}          & modifié & Aperçu du type 9 \\
\code{src/\allowbreak services/\allowbreak notificationService.js}  & modifié & \code{notifyTripAlert} \\
\code{src/\allowbreak notifications/\allowbreak notificationContract.js} & modifié & \code{buildTripPayload} \\
\code{src/services/jobQueue.js}             & modifié & \code{JOB\_WORKER\_ENABLED} \\
\code{src/\allowbreak controllers/\allowbreak contactListController.js} & modifié & Protection de \code{deleteList},
                                                            exclusion des bloqueurs \\
\code{src/\allowbreak services/\allowbreak accountDeletionService.js} & modifié & Trajets à la suppression \\
\bottomrule
\end{tabularx}
\end{center}

\section{Mobile}

\begin{center}
\begin{tabularx}{\linewidth}{@{}p{6.4cm}p{1.6cm}X@{}}
\toprule
\thd{Fichier} & \thd{Nature} & \thd{Objet} \\
\midrule
\code{lib/\allowbreak core/\allowbreak db/\allowbreak app\_database.dart}       & modifié & Trois tables, \code{schemaVersion} 24 \\
\code{lib/\allowbreak core/\allowbreak services/\allowbreak trip\_repository.dart} & nouveau & REST et Drift \\
\code{lib/\allowbreak core/\allowbreak services/\allowbreak trip\_\allowbreak socket\_\allowbreak service.dart} & nouveau & Événements temps réel \\
\code{lib/\allowbreak core/\allowbreak services/\allowbreak trip\_session\_guard.dart} & nouveau & Cycle de vie, régimes GPS \\
\code{…/\allowbreak kotlin/\allowbreak …/\allowbreak TripLocation\allowbreak ForegroundService.kt} & nouveau & Service en avant-plan \\
\code{…/\allowbreak kotlin/\allowbreak …/\allowbreak TripLocationBridge.kt}     & nouveau & Canal de méthodes \\
\code{lib/screens/trips/}                   & nouveau & Hub, composeur, suivi \\
\code{lib/\allowbreak widgets/\allowbreak chat/\allowbreak trip\_message\_card.dart} & nouveau & La carte de type 9 \\
\code{lib/\allowbreak widgets/\allowbreak trips/\allowbreak trip\_banner.dart}  & nouveau & Bandeau persistant \\
\code{lib/\allowbreak screens/\allowbreak chats/\allowbreak chat/\allowbreak chat\_bubbles.dart} & modifié & \code{case 9} et repli \\
\code{lib/\allowbreak screens/\allowbreak chats/\allowbreak chat/\allowbreak chat\_actions.dart} & modifié & Menu position à deux entrées \\
\code{lib/\allowbreak screens/\allowbreak profile/\allowbreak profile\_screen.dart} & modifié & Entrée « Trajets de confiance » \\
\code{lib/\allowbreak core/\allowbreak services/\allowbreak push\_service.dart} & modifié & Canal d'alerte, lien profond \\
\code{android/\allowbreak app/\allowbreak src/\allowbreak main/\allowbreak AndroidManifest.xml} & modifié & Deux permissions, un service \\
\code{ios/Runner/Info.plist}                & modifié & Deux clés \\
\code{lib/l10n/app\_fr.arb}, \code{app\_en.arb} & modifié & Libellés \\
\bottomrule
\end{tabularx}
\end{center}

\section{Administration}

\begin{center}
\begin{tabularx}{\linewidth}{@{}p{6.4cm}p{1.6cm}X@{}}
\toprule
\thd{Fichier} & \thd{Nature} & \thd{Objet} \\
\midrule
\code{app/\allowbreak (dashboard)/\allowbreak trips/\allowbreak page.tsx}       & nouveau & Compteurs agrégés, sans identité \\
\code{src/\allowbreak routes/\allowbreak admin/\allowbreak trips.js} \emph{(backend)} & nouveau & Une seule route de statistiques \\
\bottomrule
\end{tabularx}
\end{center}

\chapter{Vérification}

\section{Backend}

Appliquer \code{051\_trips.sql}, vérifier que le worker de jobs tourne, redémarrer, puis~:

\begin{enumerate}[leftmargin=1.6em]
  \item démarrer un trajet $\rightarrow$ 201~; en démarrer un second $\rightarrow$ 409
        \code{TRIP\_ALREADY\_ACTIVE}~;
  \item démarrer avec une liste Confiance vide $\rightarrow$ 409
        \code{TRUST\_LIST\_EMPTY}~; démarrer avec cinq membres $\rightarrow$ cinq lignes
        dans \code{trip\_watcher}, et cinq notifications~;
  \item envoyer un tableau \code{watchers} dans le corps de la requête~: il doit être
        \textbf{ignoré}, pas honoré --- l'audience se lit en base, jamais dans la
        requête~;
  \item poser un ETA à $+2$ min, ne rien faire~: vérifier le passage en
        \code{awaiting\_confirm}, les trois relances, puis l'alerte à $+10$ min~;
  \item prolonger de 15 min et vérifier que \textbf{l'ancienne échéance n'a pas survécu}
        dans \code{job\_queue} --- c'est le piège de \code{dedupe\_key}~;
  \item émettre une position depuis un second appareil $\rightarrow$ 403
        \code{DEVICE\_NOT\_OWNER}~; \code{trip:claim\_device}, puis vérifier que le premier
        s'est arrêté~;
  \item faire bloquer le propriétaire par un destinataire pendant le trajet~: il ne doit
        plus rien recevoir, et le propriétaire ne doit pas apprendre le blocage~;
  \item révoquer le dernier destinataire, et vérifier la clôture en
        \code{closed\_unwatched} après 5 min~;
  \item redémarrer le serveur en plein trajet~: l'échéance doit toujours partir.
\end{enumerate}

\section{Client}

\code{flutter pub get} $\rightarrow$
\code{dart run build\_runner build -{}-delete-conflicting-outputs} $\rightarrow$
\code{flutter run}.

\begin{enumerate}[leftmargin=1.6em]
  \item \textbf{Migration 23 $\rightarrow$ 24} sur un appareil déjà installé~: la mise à jour
        ne doit \textbf{pas} vider les données existantes.
  \item La carte de type 9 change d'état sans remonter la conversation ni créer de non-lu.
  \item \textbf{Hors ligne}~: réseau coupé pendant 10 minutes, puis rétabli --- les points
        tamponnés remontent avec leurs horodatages réels, sans doublon.
  \item Application balayée depuis les récents~: le service se relance et l'émission reprend.
  \item Notification persistante~: elle nomme les destinataires, et « Arrêter » fonctionne
        en un appui.
  \item Client ancien~: la carte de type 9 affiche le repli, pas du JSON.
\end{enumerate}

\section{Recette terrain --- ce qui ne se teste que dehors}

\begin{enumerate}[leftmargin=1.6em]
  \item Marcher jusqu'à la destination et vérifier que l'arrivée est détectée --- puis
        \textbf{qu'elle ne clôt rien}.
  \item Passer à 100 m de la destination sans s'arrêter~: aucune détection ne doit se
        produire (hystérésis et portillon de vitesse).
  \item Traverser un tunnel ou un parking souterrain~: le cercle doit lire « position
        indisponible », \textbf{pas} une alerte.
  \item Un trajet en voiture de 45 minutes, batterie mesurée avant et après.
  \item Téléphone en veille, écran éteint, 20 minutes~: le flux ne doit pas s'arrêter sur
        Android~; sur iOS, constater le comportement réel et l'écrire.
\end{enumerate}

\section{Points ouverts}

\begin{enumerate}[leftmargin=1.6em]
  \item Rétention de 12 mois du fait \code{trip} --- alignée sur la politique générale de
        conservation du produit~?
  \item \textbf{Aucune sous-sélection de l'audience} --- décision prise, et c'est ce qui
        rend le départ instantané. Le point à surveiller après la mise en service~: les
        situations où l'on voudrait partager à trois sur cinq (rendez-vous médical,
        entretien d'embauche). La réponse actuelle est de sortir la personne de son cercle
        avant de partir, ce qui est un geste lourd. Si l'usage montre que le cas est
        fréquent, l'ajout d'une sous-sélection ne coûterait rien au modèle de données~:
        \code{trip\_watcher} est déjà une table d'appartenance par trajet.
  \item Le contournement de « Ne pas déranger » doit-il être désactivable par le
        destinataire~? (Recommandation~: non, mais l'arbitrage engage le produit.)
  \item Les destinataires se voient-ils entre eux pendant un trajet --- le nombre seulement,
        ou les identités~? (Recommandation~: le nombre. Les identités exposent le carnet
        d'adresses, mais faciliteraient la coordination en cas de SOS.)
  \item Qui porte l'AIPD, et bloque-t-elle la mise en production~? (Recommandation~: oui.)
  \item Le verrou de suppression de 30 jours est-il tenable dans toutes les juridictions
        visées~?
  \item \code{BROADCAST\_WORKER\_ENABLED} $\rightarrow$ \code{JOB\_WORKER\_ENABLED}~:
        renommage acceptable, ou faut-il un second drapeau~?
  \item Quel fournisseur de tuiles cartographiques, et à quel coût par trajet~?
  \item Extensions volontairement hors périmètre, possibles sans changer les tables~:
        déclenchement matériel du SOS, appel automatique des secours, partage à un
        non-contact par lien.
\end{enumerate}
